\documentclass[article, a4paper, twoside]{universal}

\setshowlvl{0}
\begin{document}
\confighead{}{}{}
\printhead{}{}{1}


\sct{Modularity for abelian surfaces}

Ref:~\cite{BCGP2025}.

\ssc{Solid representations and classicality}

\begin{stp}
    Let $(G,X)$ be a Shimura datum of Hodge type with cocharacter $\mu$ of $G_{\bC}$, let $E$ be a finite extension of $\bQ_p$ equipped with a fixed map from the reflex field $E(G,X)$. Let $P$ be the parabolic corresponding to $\mu$ with Levi $M$, let ${}^{M}W$ be the Kostant representatives of $M$, and $w_{0}^{M}\in{}^{M}W$ the representative of maximal length.

    Let $\Fl:=P\bs G$ be the partial flag variety, $\Sh_{K_{p}K^{p}}^{\T{tor}}\ar\Spa(\bC_{p},\cO_{\bC_{p}})$ the toroidal compactification of the Shimura variety of level $K_{p}K^{p}$ for a fixed choice of cone decomposition $\Sgm$, and write $\Sh_{K^{p}}^{\T{tor}}:=\lim_{K_{p}}\Sh_{K_{p}K^{p}}^{\T{tor}}$, it is perfectoid with Hodge-Tate period map $\pi_{\T{HT}}:\Sh_{K^{p}}^{\T{tor}}\ar\Fl$, one has the smooth and locally analytic subsheaves $\cO_{\Sh_{K^{p}}^{\T{tor}}}^{\T{sm}}\sbe\cO_{\Sh_{K^{p}}^{\T{tor}}}^{\T{la}}\sbe\cO_{\Sh_{K^{p}}^{\T{tor}}}$ and the ideal sheaves $\cI_{\Sh_{K^{p}}^{\T{tor}}}^{\T{sm}}\sbe\cI_{\Sh_{K^{p}}^{\T{tor}}}^{\T{la}}\sbe\cI_{\Sh_{K^{p}}^{\T{tor}}}$ of the boundary divisor.

    Pick a maximal torus, Borel and standard parabolic $T\sbs B\sbe Q\sbs G$, for $\blt\in\{B,Q\}$, write $j_{w,\blt}:C_{w,\blt}:=P\bs Pw\blt\ahr \Fl$, $C_{w,\blt}^{\dag}:=(C_{w,\blt},j_{w,\blt}^{-1}\cO_{\Fl})$ admits an action by $\kg\rtm \blt$, write as well $j_{w,\blt}:\pi_{\T{HT}}^{-1}(C_{w,\blt})\ar\Sh_{K^{p}}^{\T{tor}}$, for a sheaf $\cF$ of solid $E$-modules on $\Sh_{K^{p}}^{\T{tor}}$, the local cohomologies $R\Gma_{w}(\Sh_{K^{p}}^{\T{tor}},\cF):=R\Gma(\Sh_{K^{p}}^{\T{tor}},j_{w,\blt,!}\cF|_{\pi_{\T{HT}}^{-1}(C_{w,\blt})})$ admit the Cousin spectral sequence (\cite[4.6.13]{BCGP2025}),
    \[
        E_{1}^{i,j}:=\Op_{w\in{}^{M}W, \ell(w)=i}\HH{w}{i+j}{(\Sh_{K^{p}}^{\T{tor}},\cF)}\aoR\HH{}{i+j}{(\Sh_{K^{p}}^{\T{tor}},\cF)}.
    \]


    The higher Coleman functors $\HC_{w,\lda},\HC_{w,\lda,\T{cusp}}:\cO(\km_{w},\kb\cap\km_{w})_{\lda\T{-alg}}^{\T{op}}\ar D(\Mod_{B(\bQ_{p})}^{\lda\T{-sm}}(E))$ and their finite slope part $\HC_{w,\lda}^{\T{fs}},\HC_{w,\lda,\T{cusp}}^{\T{fs}}:\cO(\km_{w},\kb\cap\km_{w})_{\lda\T{-alg}}^{\T{op}}\ar D(\Mod_{B(\bQ_{p})}^{\lda\T{-sm}}(E))\ar D(\Mod_{T(\bQ_{p})}^{\lda\T{-sm}}(E))$ are given by
    \begin{align*}
      \HC_{w,\lda}(M)&:=R\Gma_{w}(\Sh_{K^{p}}^{\T{tor}},\VB^{0}\cc\HCS_{w,\lda}(M)),& \HC_{w,\lda,\T{cusp}}(M)&:=R\Gma_{w}(\Sh_{K^{p}}^{\T{tor}},\VB^{0}\cc\HCS_{w,\lda}(M)\ot\cI_{\Sh_{K^{p}}^{\T{tor}}}^{\T{sm}}),\\
      \HC_{w,\lda}^{\T{fs}}(M)&:=R\Gma_{w}(\Sh_{K^{p}}^{\T{tor}},\VB^{0}\cc\HCS_{w,\lda}(M))^{\T{fs}},& \HC_{w,\lda,\T{cusp}}^{\T{fs}}(M)&:=R\Gma_{w}(\Sh_{K^{p}}^{\T{tor}},\VB^{0}\cc\HCS_{w,\lda}(M)\ot\cI_{\Sh_{K^{p}}^{\T{tor}}}^{\T{sm}})^{\T{fs}},
    \end{align*}
    they have bounded cohomological amplitude and admits a slope estimate (\cite[4.6.45, 4.6.58, 4.6.60]{BCGP2025}),
    \begin{itm}
        \item $\HC_{w,\lda},\HC_{w,\lda}^{\T{fs}}$ has amplitude $[\ell(w),d]$, and $\HC_{w,\lda,\T{cusp}},\HC_{w,\lda,\T{cusp}}^{\T{fs}}$ has amplitude $[0,\ell(w)]$.
        \item If $M\in\cO(\km_{w},\kb\cap\km_{w})_{\lda\T{-alg}}$ is generated by a highest weight vector of weight $\nu$, assume $(G,X)$ is proper or in the Siegel case, then the slopes appearing in $\HC_{w,\lda}^{\T{fs}}(M)$ and $\HC_{w,\lda,\T{cusp}}^{\T{fs}}(M)$ are $\lqs\lda-\nu+w\xI w_{0,M}\rho+\rho$.
    \end{itm}
\end{stp}


\begin{thm}[{\cite[4.7.1, 4.7.2, 4.7.5]{BCGP2025}}]
    For $\lda\in X^{*}(T)_{E}$, the functors of completed cohomology $\CC_{\lda},\CC_{\lda,\T{cusp}}:\cO(\kg,\kb)_{\lda\T{-alg}}\ar D(\Mod_{B(\bQ_{p})}^{\lda\T{-sm}}(E))$ are defined for $M\in\cO(\kg,\kb)_{\lda\T{-alg}}$ via
    \[
        \CC_{\lda}(M):=R\Hom{\kg}{}{(M,R\Gma(\Sh_{K^{p}},\bQ_{p})^{\T{la}})},\quad \CC_{\lda,\T{cusp}}(M):=R\Hom{\kg}{}{(M,R\Gma_{c}(\Sh_{K^{p}},\bQ_{p})^{\T{la}})},
    \]
    they satisfy that (1)
    \[
        \CC_{\lda}(M)\ot\bC_{p}=R\Gma(\Sh_{K^{p}}^{\T{tor}},\VB^{\T{red}}(\Loc(M))),\quad \CC_{\lda,\T{cusp}}(M)\ot\bC_{p}=R\Gma(\Sh_{K^{p}}^{\T{tor}},\VB^{\T{red}}(\Loc(M))\ot\cI_{\Sh_{K^{p}}^{\T{tor}}}^{\T{sm}}),
    \]
    (2) if $\lda$ is non-Liouville, then there are converging spectral sequences
    \begin{align*}
      E_{1}^{i,j}:=\Op_{w\in{}^{M}W, \ell(w)=i}\HH{}{i+j}{(\HC_{\oga,\lda}(M\ot_{\ku_{\kp_{w}}}^{L}E))}&\aoR\HH{}{i+j}{(\CC_{\lda}(M))}\ot\bC_{p}.\\
      E_{1}^{i,j}:=\Op_{w\in{}^{M}W, \ell(w)=i}\HH{}{i+j}{(\HC_{\oga,\lda,\T{cusp}}(M\ot_{\ku_{\kp_{w}}}^{L}E))}&\aoR\HH{}{i+j}{(\CC_{\lda,\T{cusp}}(M))}\ot\bC_{p}.
    \end{align*}
    (3) if the Shimura variety is proper or in the Siegel case, then $\CC_{\lda}^{\T{fs}}(M(\kg)_{\lda})$ and $\CC_{\lda,\T{cusp}}^{\T{fs}}(M(\kg)_{\lda})$ has slope $\gqs0$, there is converging spectral sequence
    \[
        E_{1}^{i,j}:=\Op_{w\in{}^{M}W, \ell(w)=i}\HH{}{2i+j-d}{(\HC_{w,\lda}^{=0}(M(\km_{w})_{\lda+w^{-1}w_{0,M}\rho+\rho}))}\aoR\HH{}{i+j}{(\CC_{\lda}^{=0}(M(\kg)_{\lda}))}\ot\bC_{p}.
    \]
\end{thm}

\begin{thm}[{\cite[4.5.20]{BCGP2025}}]\label{thm:LBproperty}
    Let $U_{\Fl}$ be a quasi-compact open subset of $\Fl$ over $\Spa(\bC_{p},\cO_{\bC_{p}})$, let $U_{K^{p},\Sgm}:=\pi_{\T{HT},\Sgm}\xI(U_{\Fl})$, define the functor $\VB_{\Sgm}^{0}:\LB_{\kg}(U_{\Fl})\ar\Mod(\cO_{U_{K^{p},\Sgm}}^{\T{sm}})$ via $\cF\amt\ilim_{K_{p}}(\pi_{\T{HT},\Sgm}\xI\cF\ot_{\pi_{\T{HT},\Sgm}\xI\cO_{\Fl}}\cO_{U_{K^{p},\Sgm}})^{K_{p}}$, it satisfies the following properties,
    \begin{itm}
        \item The diagram is commutative, the horizontal arrows are exact.
        \[
            \begin{tikzcd}
                \LB_{(\kg,G)}(\Fl)^{\ku_{P}^{0}}\ar[r, "{\VB^{0}}"]\ar[d, "{\Res}"] & \Mod_{G(\bQ_{p})}(\cO_{\Sh_{K^{p}}^{\T{tor}}}^{\T{sm}})\ar[d, "{\Res}"] \\
                \LB_{(\kg,G)}(C_{w}^{\dag})^{\ku_{P}^{0}}\ar[r, "{\VB^{0}}"] & \Mod_{B(\bQ_{p})}(\cO_{\pi_{\T{HT}}^{-1}C_{w}^{\dag}}^{\T{sm}})
            \end{tikzcd}
        \]
        \item One has $\VB^{0}(\cC^{\T{la}})=\cO_{\Sh_{K^{p}}^{\T{tor}}}^{\T{la}}$; if $\cF$ arises rationally, then $\Tta_{\T{hor}}(\mu)$ is an arithmetic Sen operator on $\VB^{0}(\cF)$.
        \item If $\cF\in\LB_{(\kg,G)}(\Fl)^{\ku_{P}^{0}}$, one has $\VB^{0}(\cF)\ot_{\cO_{\Sh_{K^{p}}^{\T{tor}}}^{\T{sm}}}\cO_{\Sh_{K^{p}}^{\T{tor}}}=\pi_{\T{HT}}^{-1}\cF\ot_{\pi_{\T{HT}}^{-1}\cO_{\Fl}}\cO_{\Sh_{K^{p}}^{\T{tor}}}$.
        \item If $\cF\in\LB_{(\kg,G)}(C_{w}^{\dag})^{\ku_{P}^{0}}$, one has $\VB^{0}(\cF)\ot_{\cO_{\pi_{\T{HT}}^{-1}C_{w}^{\dag}}^{\T{sm}}}\cO_{\pi_{\T{HT}}^{-1}C_{w}^{\dag}}=\pi_{\T{HT}}^{-1}\cF\ot_{\pi_{\T{HT}}^{-1}\cO_{C_{w}^{\dag}}}\cO_{\pi_{\T{HT}}^{-1}C_{w}^{\dag}}$.
    \end{itm}
\end{thm}


\spg{Case of $\GSp_{4}$}

\begin{stp}
    Now take $G=\GSp_{4}$, let $P$ be the Siegel parabolic attached to $\mu=(-1/2,-1/2;1/2)\in X_{*}(T)_{E}$, the positive simple roots are $\afa=(1,-1;0)$ and $\bta=(-2,0;0)$, let $Q$ be the Klingen parabolic attached to $\bta$, the Kostant representatives are ${}^{M}W=\{{}^{0}w:=\Id,{}^{1}w:=s_{\bta},{}^{2}w:=s_{\bta}s_{\afa},{}^{3}w:=s_{\bta}s_{\afa}s_{\bta}\}$, the maximal length representative is $w_{0}^{M}={}^{3}w$, one has the stratification $C_{{}^{3}w,Q}=C_{{}^{3}w,B}\cup C_{{}^{2}w,B}$.
\end{stp}

\begin{thm}[{\cite[3.6.9]{BCGP2025}}]\label{thm:equalExt}
    Consider $j:C_{{}^{3}w,B}^{\dag}\ahr C_{{}^{3}w,Q}^{\dag}$, take $\eta=(0,-1;1)$ and $\lda=(1,1;w)$, one has two short exact sequences of $(\kg,B)$-equivariant sheaves
    \[
        \begin{tikzcd}
            0\ar[r]& j\zE\HCS_{{}^{3}w,\eta}(M(\km_{{}^{3}w})_{\lda})\ar[r]&\HCS_{Q,{}^{3}w,\eta}(M(\km_{{}^{3}w})_{\lda})\ar[r]&\HCS_{{}^{2}w,\eta}(M(\km_{{}^{2}w})_{s_{\bta}\lda})\ar[r]& 0 \\
            0\ar[r]& j\zE\HCS_{{}^{3}w,\lda}(M(\km_{{}^{3}w})_{\lda})\ar[r]&\Ext{\kb,*_{2};\mu}{1}{((\lda,\mu_{0}),\cC^{\T{la}}|_{C_{{}^{3}w,Q}^{\dag}})}\ar[r]&\HCS_{{}^{2}w,\lda}(M(\km_{{}^{2}w})_{s_{\bta}\lda})\ar[r]& 0
        \end{tikzcd}
    \]
    they differ only by a twist of the $B$-action by character $\lda-\eta$ and multiplication by a scalar in $E\xT$.
\end{thm}

\begin{thm}[{\cite[4.9.2]{BCGP2025}}]
    The following are complement to genral cohomological amplitude results.
    
    For all $\kpa\in X^{*}(T)_{E}$, one has
    \begin{align*}
      \HH{w}{0}{(\Sh_{K^{p}}^{\T{tor}},\oga_{w}^{\dag,\kpa}(-D))^{\T{fs}}}=\HH{w}{0}{(\Sh_{K^{p}}^{\T{tor}},\oga_{w}^{\dag,\kpa})^{\T{fs}}}=0,&\quad \T{for $w\neq \Id$} \\
      \HH{w}{3}{(\Sh_{K^{p}}^{\T{tor}},\oga_{w}^{\dag,\kpa}(-D))^{\T{fs}}}=\HH{w}{3}{(\Sh_{K^{p}}^{\T{tor}},\oga_{w}^{\dag,\kpa})^{\T{fs}}}=0,&\quad \T{for $w\neq w_{0}^{M}$}
    \end{align*}

    For all $\kpa\in X^{*}(T)^{M,+}$, one has
    \begin{align*}
      \HH{w}{0}{(\Sh_{K^{p}}^{\T{tor}},\oga_{w}^{\kpa,\T{sm}}(-D))^{\T{fs}}}=\HH{w}{0}{(\Sh_{K^{p}}^{\T{tor}},\oga_{w}^{\kpa,\T{sm}})^{\T{fs}}}=0,&\quad \T{for $w\neq \Id$} \\
      \HH{w}{3}{(\Sh_{K^{p}}^{\T{tor}},\oga_{w}^{\kpa,\T{sm}}(-D))^{\T{fs}}}=\HH{w}{3}{(\Sh_{K^{p}}^{\T{tor}},\oga_{w}^{\kpa,\T{sm}})^{\T{fs}}}=0,&\quad \T{for $w\neq w_{0}^{M}$}
    \end{align*}
\end{thm}

\begin{thm}[{\cite[4.9.9]{BCGP2025}}]
    For any irreducible representation $\oL{\rho}:\Gal_{\bQ}\ar\GSp_{4}(\oL{\bF}_{p})$, the localization $V:=R\Hom{\kb}{}{(\lda,R\Gma(\Sh_{K^{p}}^{\T{tor}},\oL{\bQ}_{p})^{\T{la}})_{\km_{\oL{\rho}}}^{\T{ord}}}=\CC_{\lda}(M(\kg)_{\lda})_{\km_{\oL{\rho}}}^{\T{ord}}$ is concentrated in degree $3$, $V_{\bC_{p}}$ admits a $\Gal_{\bQ_{p}}\tms\bT_{K^{p}}\tms T(\bQ_{p})$-equivariant filtration such that the graded pieces are given by
    \[
        \Gr^{i}V_{\bC_{p}}=\HH{{}^{i}w}{i}{(\Sh_{K^{p}}^{\T{tor}},\oga_{{}^{i}w}^{\dag,-w_{0,M}({}^{i}w\cdot\lda+2\rho^{M})})}\ot\bC_{p}(-{}^{i}w^{-1}w_{0,M}\rho-\rho)_{\km_{\oL{\rho}}}^{\T{ord}},\quad i=0,1,2,3,
    \]
    where the Sen operator acts as scalar via
    \[
        (-6+\lda_{1}+\lda_{2}-w)/2,(-4-\lda_{1}+\lda_{2}-w)/2,(-2+\lda_{1}-\lda_{2}-w)/2,(-\lda_{1}-\lda_{2}-w)/2,\quad i=0,1,2,3.
    \]
\end{thm}

\begin{thm}[{\cite[4.10.3]{BCGP2025}}]\label{thm:computeQ}
    Take $\lda=(1,1;w)$,
    consider $j:\pi_{\T{HT}}^{-1}(C_{{}^{3}w,B})\ahr\pi_{\T{HT}}^{-1}(C_{{}^{3}w,Q})$, there is a corresponding extension over $C_{{}^{3}w,Q}$
    \[
        \begin{tikzcd}
            0\ar[r] & j_{!}\oga^{(1,1;-w),\T{sm}}|_{\pi_{\T{HT}}^{-1}(C_{{}^{3}w,B})}\ar[r] & \oga^{(1,1;-w),\T{sm}}|_{\pi_{\T{HT}}^{-1}(C_{{}^{3}w,Q})}\ar[r] & \oga^{(1,1;-w),\T{sm}}|_{\pi_{\T{HT}}^{-1}(C_{{}^{2}w,B})}\ar[r] & 0
        \end{tikzcd}
    \]
    such that $R\Gma_{c}(\pi_{\T{HT}}^{-1}(C_{{}^{3}w,Q}),\oga^{(1,1;-w),\T{sm}})^{\T{ord}}\ar R\Gma(\Sh_{K^{p}}^{\T{tor}},\oga^{(1,1;-w),\T{sm}})^{\T{ord}}$ is a quasi-isomorphism computed by the induced complex $\HH{{}^{2}w}{2}{(\Sh_{K^{p}}^{\T{tor}},\oga^{(1,1;-w),\T{sm}})^{\T{ord}}}\axr{\Cou}\HH{{}^{3}w}{3}{(\Sh_{K^{p}}^{\T{tor}},\oga^{(1,1;-w),\T{sm}})^{\T{ord}}}$.
\end{thm}

\begin{thm}[{C=S, \cite[4.10.12]{BCGP2025}}]\label{thm:C=S}
    The Cousin and Sen maps agree up to a non-zero scalar,
    \[
        \Cou,\Sen:\HH{{}^{2}w}{2}{(\Sh_{K^{p}}^{\T{tor}},\oga_{{}^{2}w}^{\dag,(1,1;-w)})}\ot\bC_{p}(2,0;0)_{\km_{\oL{\rho}}}^{\T{ord}}\ar\HH{{}^{3}w}{3}{(\Sh_{K^{p}}^{\T{tor}},\oga_{{}^{3}w}^{\dag,(1,1;-w)})_{\km_{\oL{\rho}}}^{\T{ord}}}.
    \]
\end{thm}
\begin{prf}[Sketch of the proof]
    This follows from \cref{thm:LBproperty} and \cref{thm:equalExt}.
\end{prf}

\begin{thm}[{Classicality, \cite[4.12.4]{BCGP2025}}]\label{thm:classicality}
    Let $f\in\HH{{}^{0}w}{0}{(\Sh_{K^{p}}^{\T{tor}},\oga^{(2,2;-w),\T{sm}})^{\T{ord}}}$ with maximal ideal $\km_{f}$ and associated Galois representation $\rho_{f}:\Gal_{\bQ}\ar\GSp_{4}(\oL{\bQ}_{p})$, if (1) $\oL{\rho}_{f}$ is irreducible, and (2) $\rho_{f}$ has large image, and is de Rham at $p$, and (3) there exists $n\in\bZ$ such that $\Dim_{\bC_{p}}\Gr^{i}V_{\bC_{p}}[\km_{f}]=n,0\lqs i\lqs3$ and $\Dim_{\oL{\bQ}_{p}}V[\km_{f}]=4n$, where $V:=\HH{}{3}{(R\Hom{\kb}{}{(\lda,R\Gma(\Sh_{K^{p}}^{\T{tor}},\oL{\bQ}_{p})^{\T{la}})_{\km_{\oL{\rho}_{f}}}^{\T{ord}}})}$, then $f$ is a classical modular form.
\end{thm}

\begin{prf}[Sketch of the proof]
    This follows from \cref{thm:C=S} and \cref{thm:computeQ}.
\end{prf}

\ssc{Multiplicity one and modularity criterion}

\begin{thm}[{Modularity criterion, \cite[7.5.11]{BCGP2025}}]
    $\rho:\Gal_{\bQ}\ar\GSp_{4}(\cO_{E})$ is modular if
    \begin{itm}
        \item (1) $\rho$ has large image, is unramified outside finitely many primes and pure, (2) $\nu\cc\rho=\eps^{-1}$, (3) $\rho(\Gal_{\bQ(\zta_{p^{\ift}})})$ is integrally enormous, (4) $\rho|_{\Gal_{\bQ_{p}}}$ is ordinary, semistable of weight $2$ and $p$-distinguished.
        \item If $p=2$, and there is a solvable CM extension $F/\bQ$ and an ordinary RACSDC automorphic representation $\pi$ of $\GL_{4}/F$ such that (1) $\oL{r}_{\pi,2}\cong\oL{\rho}|_{\Gal_{F}}$, (2) $\oL{\rho}(\Gal_{F})$ is nearly adequate and contains a regular semi-simple element, (3) $\oL{\rho}(\Gal_{\bQ(i)})=\oL{\rho}(\Gal_{\bQ})$, (4) there is an infinite place $v$ of $F^{+}$ such that $(\oL{\rho}|_{\Gal_{F}},1)$ is strongly residually odd at $v$.
        \item If $p>2$, and there is an ordinary cuspidal automorphic representation $\pi$ of $\GSp_{4}/\bQ$ with central character $\sP{\cdot}^{2}$ such that (1) $\oL{r}_{\pi,p}\cong\oL{\rho}$, (2) $\oL{\rho}$ is $\GSp_{4}$-reasonable and tidy, (3) $\oL{\rho}(\Gal_{\bQ})\ssm\Sp_{4}(\bF_{p})$ contains a regular semisimple element, (4) there is a compatible choice of $p$-stabilization of $\pi_{p}$ and $\rho|_{\Gal_{\bQ_{p}}}$ such that $r_{\pi,p}|_{\Gal_{\bQ_{p}}}$ and $\rho|_{\Gal_{\bQ_{p}}}$ lie on the same unique component of $\Spc R_{p}^{\trg}$.
    \end{itm}
\end{thm}

\begin{prf}[Sketch of the proof]
    By higher Hida theory \cref{thm:highHida} and multiplicity one \cref{thm:mul1}, the conditions required by the classicality result \cref{thm:classicality} will be satisfied, 
\end{prf}

\begin{thm}[{\cite[7.3.1]{BCGP2025}}]\label{thm:highHida}
    For each $w\in{}^{M}W$ and a neat tame level $K^{p}$, one has perfect complexes of $\bZ_{p}\dS{T(\bZ_{p})}$-modules denoted by $M_{w}^{\blt}$ and $M_{w,\T{cusp}}^{\blt}$, define
    \[
        M_{w}:=\begin{cases}
          \HH{}{\ell(w)}{(M_{w,\T{cusp}}^{\blt})}\ot_{\bZ_{p}}\cO_{E} & \ell(w)=0,1 \\
          \HH{}{\ell(w)}{(M_{w}^{\blt})}\ot_{\bZ_{p}}\cO_{E} & \ell(w)=2,3
        \end{cases}
    \]
    they satisfy the following properties,
    \begin{itm}
        \item $M_{w}$ are finite projective $\cO_{E}\dS{T(\bZ_{p})}$-modules, and there is a perfect duality pairing
        \[
            M_{w}\ot_{\cO_{E}\dS{T(\bZ_{p})}}M_{w_{0,M}ww_{0}}\ar\cO_{E}\dS{T(\bZ_{p})},
        \]
        interpolating the classical Serre duality. For normal subgroups $K_{1}^{p}\sbs K_{2}^{p}$, $M_{w,K_{1}^{p}}$ are finite projective $\cO_{E}\dS{T(\bZ_{p})}[K_{2}^{p}/K_{1}^{p}]$-modules. The pullback and trace induces isomorphisms
        \[
            M_{w,K_{2}^{p}}\cong M_{w,K_{1}^{p}}^{K_{2}^{p}/K_{1}^{p}},\quad (M_{w,K_{1}^{p}})_{K_{2}^{p}/K_{1}^{p}}\cong M_{w,K_{2}^{p}}.
        \]
        \item for any $\lda\in X^{*}(T)^{+}$ with $\kpa:=-w_{0,M}w(\lda+\rho)-\rho$, there are Hecke-equivariant isomorphisms
        \[
            M_{w}\ot_{\cO_{E}\dS{T(\bZ_{p})},\lda}E(-\lda)=\begin{cases}
              \HH{}{\ell(w)}{(\Sh_{\Iw(p)K^{p}}^{\T{tor}},\oga^{\kpa}(-D))^{\T{ord}}} & \ell(w)=0,1\\
              \HH{}{\ell(w)}{(\Sh_{\Iw(p)K^{p}}^{\T{tor}},\oga^{\kpa})^{\T{ord}}} & \ell(w)=2,3
            \end{cases}
        \]
        \item for any $\lda\in X^{*}(T)$ with $\kpa:=-w_{0,M}w(\lda+\rho)-\rho\in X^{*}(T)^{M,+}$, there are Hecke-equivariant isomorphisms
        \[
            M_{w}\ot_{\cO_{E}\dS{T(\bZ_{p})},\lda}E(-\lda)=\begin{cases}
              \HH{w}{\ell(w)}{(\Sh_{\Iw(p)K^{p}}^{\T{tor}},\oga^{\kpa}(-D))^{\T{ord},T(\bZ_{p})}} & \ell(w)=0,1\\
              \HH{w}{\ell(w)}{(\Sh_{\Iw(p)K^{p}}^{\T{tor}},\oga^{\kpa})^{\T{ord},T(\bZ_{p})}} & \ell(w)=2,3
            \end{cases}
        \]
    \end{itm}
\end{thm}

\begin{prf}[Sketch of the proof]
    See BP2023.
\end{prf}


\begin{thm}[{Multiplicity one, \cite[7.5.10]{BCGP2025}}]\label{thm:mul1}
    Suppose $\rho:\Gal_{\bQ}\ar\GSp_{4}(\cO_{E})$ satisfies that
    \begin{itm}
        \item (1) $\rho$ is unramified at all but finitely many primes and pure, (2) $\nu\cc\rho=\eps^{-1}$, (3) $\rho(\Gal_{\bQ(\zta_{p^{\ift}})})$ is integrally enormous, (4) $\rho|_{\Gal_{\bQ_{p}}}$ is ordinary, semistable of weight $2$ and $p$-distinguished.
        \item If $p=2$, and there is a solvable CM extension $F/\bQ$ and an ordinary RACSDC automorphic representation $\pi$ of $\GL_{4}/F$ such that (1) $\oL{r}_{\pi,2}\cong\oL{\rho}|_{\Gal_{F}}$, (2) $\oL{\rho}(\Gal_{F})$ is nearly adequate and contains a regular semisimple element, (3) $\oL{\rho}(\Gal_{\bQ(i)})=\oL{\rho}(\Gal_{\bQ})$, (4) there is an infinite place $v$ of $F^{+}$ such that $(\oL{\rho}|_{\Gal_{F}},1)$ is strongly residually odd at $v$.
        \item If $p>2$, and there is an ordinary cuspidal automorphic representation $\pi$ of $\GSp_{4}/\bQ$ with central character $\sP{\cdot}^{2}$ such that (1) $\oL{r}_{\pi,p}\cong\oL{\rho}$, (2) $\oL{\rho}$ is $\GSp_{4}$-reasonable and tidy, (3) $\oL{\rho}(\Gal_{\bQ})\ssm\Sp_{4}(\bF_{p})$ contains a regular semisimple element, (4) there is a compatible choice of $p$-stabilizations of $\pi_{p}$ and $\rho|_{\Gal_{\bQ_{p}}}$ such that $r_{\pi,p}|_{\Gal_{\bQ_{p}}}$ and $\rho|_{\Gal_{\bQ_{p}}}$ lie on the unique component of $\Spc{R_{p}^{\trg}}$.
    \end{itm}
    Then $\rho$ automatically satisfies (7.1.3), let $r$ be a neat prime of $\oL{\rho}$, $S$ be the union of $\{r\}$ and ramified primes of $\rho$, let $\cS$ be the global deformation problem $\cS:=(\oL{\rho},S,\{\Lda_{v}\}_{v\in S_{p}}\cup\{\cO_{E}\}_{v\in S\bs S_{p}},\{\cD_{v}^{\trg}\}_{v\in S_{p}}\cup\{\cD_{v}^{\wsq}\}_{v\in S\bs S_{p}})$. Take the level $K^{p}=\prd_{l}K_{l}\sbs\GSp_{4}(\bA^{\ift,p})$ defined via $K_{l}=\GSp_{4}(\bZ_{l})$ if $l\nin S$, $K_{r}=\Iw_{1}(r)$, and $K_{l}=\Par(l^{a(\rho|_{\Gal_{\bQ_{l}}})})$ if $l\in S\ssm\{r\}$. Then take $M:=\Hom{\Lda_{\GSp_{4},\bQ}}{}{((M_{K^{p}}\ot_{\cO\dS{T(\bZ_{p})}}\Lda_{\GSp_{4},\bQ})^{\sP{\cdot}^{2}},\Lda_{\GSp_{4},\bQ})_{\km,\km_{r}}}$.

    Let $\kq:=\Ker(R_{\cS}\ar\cO_{E})$ where the morphism corresponds to $\rho$, then $M_{\kq}$ is a free $(R_{\cS})_{\kq}$-module of rank $1$.
\end{thm}

\begin{prf}[Sketch of the proof]
    $M_{\kq}$ is a free $(R_{\cS})_{\kq}$-module by \cref{thm:Rfree}. If $p=2$, then $R_{\cS}^{\T{red}}[1/p]\ar\bT_{\cS}[1/p]$ is an isomorphism combining solvable descent \cref{thm:solvable} and a modularity lifting at $2$ \cref{thm:p=2-R=T}; if $p>2$, then $\Img(\Spc\bT_{\cS}[1/p]\ar\Spc R_{\cS}^{\T{red}}[1/p])$ contains an irreducible component containing $\kq$ combining solvable descent \cref{thm:solvable} and a modularity lifting at $p$ \cref{thm:p>2-R=T}. Thus $M_{\kq}$ is nonzero in either case, it is free of rank one again by a multiplicity one result \cite[7.5.8]{BCGP2025}.
\end{prf}

\begin{thm}[{\cite[7.2.3]{BCGP2025}}]\label{thm:Rfree}
    Let $F^{+}$ be a totally real number field in which $p$ splits completely and fix $\rho:\Gal_{F^{+}}\ar\GSp_{4}(\cO_{E})$ satisfying that (7.1.3),
    \begin{itm}
        \item $\rho$ is (1) unramified outside finitely many places and pure, (2) $\nu\cc\rho=\eps^{-1}$, (3) $\rho(\Gal_{F^{+}(\zta_{p^{\ift}})})$ is integrally enormous, (4) for all $v\in S_{p}$, $\rho|_{\Gal_{F_{v}}^{+}}$ is ordinary, semistable of weight $2$, pure and $p$-distinguished.
        \item $\oL{\rho}:\Gal_{F^{+}}\ar\GSp_{4}(k)$ is absolutely irreducible; if $p=2$, then $\oL{\rho}(\Gal_{F^{+}(i)})=\oL{\rho}(\Gal_{F^{+}})$; if $p>2$, then $\oL{\rho}(\Gal_{F^{+}})\ssm\Sp_{4}(\bF_{p})$ contains a regular semi-simple element.
    \end{itm}
    Let $S$ be a finite set of places of $F^{+}$ containing $S_{p}$ and ramified places of $\rho$, consider the global deformation problem $\cS:=(\oL{\rho},S,\{\Lda_{v}\}_{v\in S_{p}}\cup\{\cO_{E}\}_{v\in S\bs S_{p}},\{\cD_{v}^{\trg}\}_{v\in S_{p}}\cup\{\cD_{v}^{\wsq}\}_{v\in S\bs S_{p}})$, let $Q$ be a set of Taylor-Wiles primes, the augmented global deformation problem associated to $Q$ is
    \[
        \cS_{Q}:=(\oL{\rho},S\cup Q,\{\Lda_{v}\}_{v\in S_{p}}\cup\{\cO_{E}\}_{v\in (S\cup Q)\bs S_{p}},\{\cD_{v}^{\trg}\}_{v\in S_{p}}\cup\{\cD_{v}^{\wsq}\}_{v\in (S\cup Q)\bs S_{p}}).
    \]
    Let $M$ be a $R_{\cS}$-module, finite free as a $\Lda_{\GSp_{4},F^{+}}$-module, and for each $N\gqs1$ let $M_{N}$ be a $R_{\cS_{Q_{N}}}$-module, finite free as a $\Lda_{\GSp_{4},F^{+}}[\Dta_{N}]$-module, such that there are isomophisms of $\Lda_{\GSp_{4},F^{+}}$-modules $M_{N}\ot_{\Lda_{\GSp_{4},F^{+}}[\Dta_{N}]}\Lda_{\GSp_{4},F^{+}}\cong M$, compatible with the actions of $R_{\cS_{Q_{N}}}\ar R_{\cS}$ (7.2.1).

    Let $\kq:=\Ker(R_{\cS}\ar\cO_{E})$ where the morphism corresponds to $\rho$, then $M_{\kq}$ is a finite free $(R_{\cS})_{\kq}$-module.
\end{thm}

\begin{thm}[{\cite[5.7.14]{BCGP2025}}]\label{thm:p=2-R=T}
    Let $F$ be a CM field, $n\gqs2$, and $\rho:\Gal_{F}\ar\GL_{n}(\oL{\bQ}_{2})$ satisfying that
    \begin{itm}
        \item $F/F^{+}$ is everywhere unramified, all places of $F^{+}$ which is (1) above $2$, or (2) under a ramified place of $\pi$ split in $F$. $n[F^{+}:\bQ]\eqv0\bmod{4}$.
        \item There is an ordinary RACSDC automorphic representation $\pi$ of $\GL_{n}/F$ such that $\oL{r}_{\pi,2}\cong\oL{\rho}$.
        \item $\oL{\rho}(\Gal_{F})$ is nearly adequate and contains a regular semi-simple element.
        \item If $n$ is even, there is an infinite place $v$ of $F^{+}$ such that $(\oL{\rho},\oL{\eps}^{1-n}\dta_{F/F^{+}}^{n})$ is strongly residually odd at $v$.
    \end{itm}
    Let $T$ be any finite set of places of $F^{+}$ which split in $F$, then $R^{T,\T{ord}}$ is a finite $\Lda$-algebra. If furthermore $U_{v}$ is sufficiently deep for each $v\in T\ssm S_{2}$, then the morphism $R^{T,\T{ord}}\ar\bT^{T,\T{ord}}(U(\kl^{\ift}))_{\km}$ has nilpotent kernel.
\end{thm}

\begin{thm}[{\cite[6.3.4]{BCGP2025}}]\label{thm:p>2-R=T}
    Let $F^{+}$ be totally real in which $p>2$ splits completely, and $\pi$ an ordinary cuspidal automorphic representation for $\GSp_{4}/F^{+}$ of central character $\sP{\cdot}^{2}$ and weight $(k_{v},l_{v};2)_{v\mid\ift}$, let $\rho_{\pi,p}:\Gal_{F^{+}}\ar\GSp_{4}(\oL{\bQ}_{p})$ be the associated Galois representation, let $R$ be a finite set of places of $F^{+}$ prime to $p$ containing ramified places of $\pi$, such that the following holds (6.3.1),
    \begin{itm}
        \item $\oL{\rho}_{\pi,p}$ is $\GSp_{4}$-reasonable and tidy, choose an unramified place $w_{0}$ satisfying tidyness, let $S:=S_{p}\cup R\cup \{w_{0}\}$.
        \item For each $v\in S_{p}$, $\oL{\rho}_{\pi,p}|_{\Gal_{F_{v}^{+}}}$ is ordinary of weight $2$ with a fixed $p$-stabilization $(\oL{\afa}_{v},\oL{\bta}_{v})$ compatible with a fixed choice of ordinary $p$-stabilization of $\pi_{v}$, $\rho_{\pi,p}|_{\Gal_{F_{v}^{+}}}$ lies on a unique irreducible component of $R_{v}^{\trg}$.
        \item For each place $v\in R$ one has (1) $\oL{\rho}_{\Gal_{F_{v}^{+}}}$ is trivial, $q_{v}\eqv1\bmod{p}$, (2) $q_{v}\eqv1\bmod{9}$ if $p=3$, (3) $\pi_{v}^{\Iw(v)}\neq0$.
    \end{itm}
    Consider the global deformation problem
    \[
        \cS_{1,\pi}=(\oL{\rho}_{\pi,p},S,\{\Lda_{\GSp_{4},v}\}_{v\in S_{p}}\cup\{\cO_{E}\}_{v\in S\ssm S_{p}},\{\cD_{v}^{\trg,\pi}\}_{v\in S_{p}}\cup\{\cD_{v}^{1}\}_{v\in R}\cup\{\cD_{w_{0}}^{\wsq}\}),
    \]
    then $R_{\cS_{1,\pi}}$ is a finite $\Lda_{\GSp_{4},F^{+}}$-algebra, the morphism $R_{\cS_{1,\pi}}\ar\bT_{\cS_{1,\pi}}$ has nilpotent kernel.
\end{thm}

\begin{thm}[{Solvable descent, \cite[7.5.3]{BCGP2025}}]\label{thm:solvable}
    The following are true,
    \begin{itm}
        \item If $\Pi$ is a $\sP{\cdot}^{2}$-self dual RAC automorphic representation of $\GL_{4}/\bQ$, then $(\Pi,\sP{\cdot}^{2})$ is of symplectic type.
        \item Let $F$ be a CM field so that $F/\bQ$ is a solvable Galois extension, if a representation $\rho:\Gal_{\bQ}\ar\GSp_{4}(\oL{\bQ}_{p})$ satisfies that (1) it has multiplier $\eps^{-1}$, (2) $\rho|_{\Gal_{F}}$ is irreducible, (3) there is a RACSDC automorphic representation $\pi_{F}$ of $\GL_{4}/F$ such that $\rho|_{\Gal_{F}}\cong\rho_{\pi_{F},p}\ot\eps$, then there is a RAC automorphic representation $\pi$ of $\GSp_{4}/\bQ$ with central character $\sP{\cdot}^{2}$ such that $\rho\cong\rho_{\pi,p}$.
        \item Let $F^{+}$ be a totally real field so that $F^{+}/\bQ$ is a solvable Galois extension, if a representation $\rho:\Gal_{\bQ}\ar\GSp_{4}(\oL{\bQ}_{p})$ satisfies that (1) it has multiplier $\eps^{-1}$, (2) $\rho|_{\Gal_{F^{+}}}$ is irreducible, (3) there is a RAC automorphic representation $\pi_{F^{+}}$ of $\GSp_{4}/F^{+}$ with central character $\sP{\cdot}^{2}$ such that $\rho|_{\Gal_{F^{+}}}\cong\rho_{\pi_{F^{+}},p}$, then there is a RAC automorphic representation $\pi$ of $\GSp_{4}/\bQ$ with central character $\sP{\cdot}^{2}$ such that $\rho\cong\rho_{\pi,p}$.
    \end{itm}
\end{thm}

\ssc{Arithmetic geometry of abelian surfaces}

\begin{thm}[{$2$-$3$-switch surface, \cite[9.4.2]{BCGP2025}}]\label{thm:2-3-switch}
    If $\oL{\rho}:\Gal_{\bQ}\ar\GSp_{4}(\bF_{3})$ has similitude $\oL{\eps}^{-1}$, $\oL{\rho}|_{\Gal_{\bQ_{3}}}$ is ordinary and finite flat, $\oL{\rho}|_{\Gal_{\bQ_{2}}}$ is unramified and $\Chr_{\oL{\rho}(\Frb_{2})}(x)\neq(x^{2}\pm x +2)^{2}$, then there exists a genus two curve $X/\bQ$ with a rational Weierstrass point with $B=\Jac(X)$ such that (1) $\oL{\rho}_{B,3}\cong\oL{\rho}$, (2) $B$ has good ordinary or semistable ordinary reduction at $2$, good ordinary reduction at $3$, and is $2$-distinguished, (3) $\oL{\rho}_{B,2}:\Gal_{\bQ}\ar\GSp_{4}(\bF_{2})$ has image $S_{5}(b)$, and $\oL{\rho}_{B,2}(c)$ has conjugacy class $(**)(**)$.
\end{thm}

\begin{prf}[Sketch of the proof]
    One can use local constancy (\cite[Theorem~5.1]{Kisin1999}) and brute-force group-theoretical checking (\cite[9.2.3]{BCGP2025}) to establish the existence of $X$, and then use an approximation lemma (\cite[9.4.1]{BCGP2025}) to deduce the remaining desired properties.
\end{prf}

\begin{thm}[{Modularity transfer, \cite[9.5.1]{BCGP2025}}]\label{thm:modularity-transfer}
    Let $A,B/\bQ$ be abelian surfaces, $p$ an odd prime, if $B$ is modular and has good ordinary reduction at $p$, then $A$ is modular if $\oL{\rho}_{A,p}\cong\oL{\rho}_{B,p}$ and the following hold,
    \begin{itm}
        \item $A$ has good ordinary reduction at $p$; $\rho_{A,p}|_{\Gal_{\bQ_{p}}}$ is $p$-distinguished, $\oL{\rho_{A,p}(\Gal_{\bQ})}$ contains $\Sp_{4}$.
        \item $\oL{\rho}_{A,p}$ is $\GSp_{4}$-reasonable and tidy; $\oL{\rho}_{A,p}(\Gal_{\bQ(\zta_{p})})$ contains a regular semi-simple element; $\oL{\rho}_{A,p}(\Gal_{\bQ})\ssm\Sp_{4}(\bF_{p})$ contains a regular semi-simple element.
    \end{itm}
\end{thm}

\begin{thm}[{$2$-adic modularity, \cite[8.3.2]{BCGP2025}}]\label{thm:2-adic-modularity}
    An abelian surface $A/\bQ$ is modular if the following hold,
    \begin{itm}
        \item $A$ has good ordinary or semistable reduction at $2$, and $\rho_{A,2}|_{\Gal_{\bQ_{2}}}$ is $2$-distinguished.
        \item $A_{5}(b)\sbs\oL{\rho}_{A,2}(\Gal_{\bQ})\sbs S_{5}(b)$, $\oL{\rho}_{A,2}(c)\in A_{5}(b)$ has order $2$.
    \end{itm}
\end{thm}

\begin{thm}[{\cite[9.5.2]{BCGP2025}}]
    Let $A/\bQ$ be an abelian surface with a polarization of degree prime to $3$, $A$ is modular if the following conditions hold,
    \begin{itm}
        \item $A$ has good ordinary reduction at $3$, $\rho_{A,3}|_{\Gal_{\bQ_{3}}}$ is $3$-distinguished.
        \item $\oL{\rho}_{A,3}$ is surjective, $\oL{\rho}_{A,3}|_{\Gal_{\bQ_{2}}}$ is unramified, and $\Chr_{\oL{\rho}_{A,3}(\Frb_{2})}(x)\neq(x^{2}\pm x+2)^{2}$.
    \end{itm}
\end{thm}

\begin{prf}[Sketch of the proof]
    Indeed, by \cref{thm:2-3-switch}, there exists an abelian surface $B/\bQ$, it is modular as it satisfies the conditions of a $2$-adic modularity criterion \cref{thm:2-adic-modularity}, by applying modularity transfer \cref{thm:modularity-transfer} with $p=3$ one obtains the result.
\end{prf}

\begin{rmk}
    See \cite[6.4.3]{BCGP2025} for the crucial brute-force hard work on the $3$-reduction of an abelian surface to ensure the modularity transfer works.
\end{rmk}

\begin{cmt}    
\begin{dfn}
    Say that $\rho:\Gal_{\bQ}\ar\GSp_{4}(\oL{\bQ}_{p})$ \tbf{has large image} if either (1) $\oL{\rho(\Gal_{\bQ})}$ contains $\Sp_{4}$, or (2) $\oL{\rho(\Gal_{\bQ})}$ contains $\SL_{2}\tms\SL_{2}$, $\rho$ is absolutely irreducible and reducible on some index two subgroup.
\end{dfn}

\begin{thm}[{\cite[4.11.1, 4.11.2]{BCGP2025}}]
    Let $\rho:\Gal_{\bQ}\ar\GSp_{4}(\oL{\bQ}_{p})$ and $s:\Gal_{\bQ}\ar\GL_{n}(\oL{\bQ}_{p})$ be two continuous representations, if (1) $\rho$ has large image, and (2) for a density one set of primes $\ell$, one has $\Chr_{\rho(\Frb_{\ell})}(s(\Frb_{\ell}))=0$, then $s\cong\rho^{\op m}$ for some $m\gqs1$.
\end{thm}
\end{cmt}

\begin{cmt}

    for horizontal action, $\Tta_{\T{hor}}:Z(\km)\ar\End{\cO_{U}}{}{(\cF)}$, see 3.2.25

    See 3.2.34 for the equivalence identifying $\Tta_{\T{hor}}$ and $\Tta$ actions of $Z(\km)$.

    \red{added related BGG cat and condensed, action of $\Tta_{\T{hor}}$ and $Z(\km)$}

    For fixed $w\in{}^{M}W$, let $\cO(\km_{w},\kb\cap\km_{w})$ be the BGG category; for $\lda\in X^{*}(T)_{E}$, let $\cO(\km_{w},\kb\cap\km_{w})_{\lda\T{-alg}}$ be its subcategory with weights in $\lda+X^{*}(T)$. Let $\Mod_{B(\bQ_p)}^{\lda\T{-sm}}(E)$ be the category of locally analytic representations of $B(\bQ_{p})$ with $\kb$ acting via $\lda$.


    The higher Coleman sheaf (3.4.17) is given by $\HCS_{Q,w,\lda}:\cO(\km_{w},\kq\cap\km_{w})_{\lda\T{-alg}}^{\T{op}}\ar\LB_{(\kg,Q)}(C_{w,Q}^{\dag})^{\ku_{\kp}^{0}}$;


        \red{define For VB (4.5)}

        Consider the functor $\VB^{0}:\LB_{(\kg,G)}(\Fl)^{\ku_{P}^{0}}\ar\Mod_{G(\bQ_{p})}(\cO_{\Sh_{K^{p}}^{\T{tor}}}^{\T{sm}})$

    it is exact, could be described over an analytic covering, admits explicit expression after tensoring $\cO_{\Sh_{K^{p}}^{\T{tor}}}$, $\VB^{0}(\cC^{\T{la}})=\cO_{\Sh_{K^{p}}^{\T{tor}}}^{\T{la}}$, etc.

    \red{define for $\cC^{\T{la}}$ (3.2.37)}


        \red{define for Loc (3.5)}

        attached to $M\in\cO(\kg,\kb)$ a certain twistied $D$-module $\Loc(M)$.

        \red{(4.6.53)}

        there is a finite slope functor $(-)^{\T{fs}}:D(\Mod_{B(\bQ_{p})}^{\lda\T{-sm}}(E))\ar D(\Mod_{T(\bQ_{p})}^{\lda\T{-sm}}(E))$

        \red{for big sheaves $\oga_{w}^{\dag,\lda}$}

    (4.6.6, 4.6.8, 4.6.9) Consider the following object
    \[
        \oga_{w}^{\dag,\lda}:=\VB^{0}\cc\HCS_{w,-w^{-1}w_{0,M}\lda}(M(\km_{w})_{-w^{-1}w_{0,M}\lda})\in\Mod_{B(\bQ_{p})}^{-w^{-1}w_{0,M}\lda\T{-sm}}(\cO_{\pi_{\T{HT}}^{-1}C_{w}^{\dag}}^{\T{sm}}).
    \]

     the arithmetic Sen operator acts on $\oga_{w}^{\dag,\lda}$ via $\sA{\mu,\lda}$, and for $\lda\in X^{*}(T)^{M,+}$, there is an injective $B(\bQ_{p})$-equivariant map $\oga^{\lda,\T{sm}}|_{\pi_{\T{HT}}^{-1}(C_{w}^{\dag})}\ot E(-w^{-1}w_{0,M}\lda)\ahr\oga_{w}^{\dag,\lda}$.


        \red{for smooth sheaves $\oga_{w}^{\lda,\T{sm}}$}

        (4.6.1) let $\oga^{\lda,\T{sm}}:=\VB^{0}(\cL_{\lda})$.

        in turn, 3.2.18, for $\lda\in X^{*}(T)^{M,+}$, let $\cL_{\lda}$ be the sheaf corresponding to the highest weight representation of $M$

    Let $\Dta_{\ift}:=\bZ_{p}^{2q}$, and for each $N\gqs1$ let $\Dta_{N}$ be the finite quotient of $\Dta_{\ift}$ such that $\Dta_{\ift}\atr\Dta_{\ift}/p^{N}$ factors through $\Dta_{N}$, one has the diagram
    \[
        \begin{tikzcd}[ampersand replacement=\&]
            \Lda_{\GSp_{4},F^{+}}[\Dta_{N}]\ar[r]\ar[d] \& R_{\cS_{Q_{N}}}\ar[d]\\
            \Lda_{\GSp_{4},F^{+}}\ar[r] \& R_{\cS}
        \end{tikzcd}
    \]

    \red{Nearly adequate/weakly adequate}
    
    See 5.3.2, 5.3.3

    \red{For purity}
    
    See\cite[2.5]{BCGP2021}


    \red{for strongly residually odd} See \cite[Definition~3.3]{Thorne2017}

    \begin{dfn}[{\cite[7.5.11]{BCGP2021}}]
        Say that a subgroup $H\sbs\GSp_{4}(k)$ is \tbf{tidy} if there is $h\in H$ with $\nu(h)\neq1$ and no two eigenvalues of $h$ has ratio $\nu(h)$. Say that $\oL{\rho}:\Gal_{F}\ar\GSp_{4}(k)$ is tidy if it has tidy image.
    \end{dfn}


    \begin{dfn}[7.4.1]
        A \tbf{neat prime} for $\oL{\rho}:\Gal_{\bQ}\ar\GSp_{4}(\cO_{E})$ is a prime $r>5$ such that $r\neq p$, $r\not\eqv1\bmod{p}$ and $r\not\eqv1\bmod{4}$ if $p=2$, $\rho|_{\Gal_{\bQ_{r}}}$ is unramified, $\oL{\rho}(\Frb_{r})$ is regular semi-simple, together with a fixed ordering $(\oL{\afa}_{r,1},\oL{\afa}_{r,2},r\oL{\afa}_{r,2}^{-1},r\oL{\afa}_{r,1}^{-1})$ of the eigenvalues of $\oL{\rho}(\Frb_{r})$. If $\rho$ satisfies (7.1.3) then such $r$ exists.
    \end{dfn}


    A set of Taylor-Wiles primes of level $N$ is a finite set $Q$ of places of $F^{+}$, disjoint from $S$ such that for each $v\in Q$, one has $q_{v}\eqv1\bmod{p^{N}}$ and $\oL{\rho}(\Frb_{v})$ is regular semi-simple.


    \begin{dfn}[7.5.7]
        $R$-regular classical weight.
    \end{dfn}

    \begin{thm}[7.1.13]
        Suppose $\rho:\Gal_{F^{+}}\ar\GSp_{4}(\cO_{E})$ satisfies 7.1.3, then there exists an integer $d\gqs0$ such that for each $N\gqs1$ there is set of Taylor-Wiles primes $Q_{N}$ of level $N$, together with a morphism in $\T{CNL}_{\Lda_{\GSp_{4},F^{+}}}$, $R_{\cS,S}^{\T{loc}}\dS{X_{1},\lds,X_{g}}\ar R_{\cS_{Q_{N}}}^{S}$ satisfying certain conditions.
    \end{thm}

    \begin{dfn}[7.4.9]
        Let $\bT_{\cS}$ be the image of $\bT$ in $\End{\Lda_{\GSp_{4},\bQ}}{}{(M)}$.
    \end{dfn}

    \begin{dfn}[1.8.10]
        Let $K/\bQ_{p}$ be a finite extension, and $\rho:\Gal_{K}\ar\GSp_{4}(\oL{\bQ}_{p})$ a representation with similitude $\eps^{-1}$. Say that $\rho$ is \tbf{ordinary} if there are characters $\chi_{1},\chi_{2}$ if $\rho$ is isomorphic to an upper-triangular matrix with diagonal $\Dia(\chi_{1},\chi_{2},\eps^{-1}\chi_{2}^{-1},\eps^{-1}\chi_{1}^{-1})$, such a pair $(\chi_{1},\chi_{2})$ is said to be a \tbf{$p$-stabilization} of $\rho$; say that $\rho$ is \tbf{$p$-distinguished} if $\chi_{1},\chi_{2},\eps^{-1}\chi_{1}^{-1},\chi_{2}^{-1}$ are pairwise distinct; say that $\rho$ is semistable of weight $2$ if its upper-left representation with diagonal $(\chi_{1},\chi_{2})$ is unramified.


        (Similar for residual representations)
    \end{dfn}
    
    \begin{dfn}[9.1.2]
        Say that an abelian surface $A$ is $p$-distinguished if the corresponding representation $\rho_{A,p}:\Gal_{\bQ_{p}}\ar\GSp_{4}(\bQ_{p})$ is $p$-distinguished, say that a curve $X$ of genus $2$ is $p$-distinguished if its Jacobian $\Jac(X)$ is $p$-distinguished.
    \end{dfn}
\end{cmt}


\printref
\end{document}
